\documentclass[11pt]{article}

\usepackage[margin=1in]{geometry}
\usepackage{booktabs}
\usepackage{graphicx}
\usepackage{array}
\usepackage{longtable}
\usepackage{xcolor}
\usepackage{hyperref}

\graphicspath{{report_figures/}}

\title{Real-Ecology Data Benchmark: Experiment Results}
\author{}
\date{2026-07-05}

\begin{document}
\maketitle

\section{Summary}

This report summarizes the first audited experiment set after moving from the
dummy ecology setting to the real-ecology-data setting. The results are drawn
from the frozen run
\path{real_ecology_runs/psafe_overnight_20260705/}. The benchmark uses
continuous abundance states, noisy observations, 9 real populations, 11 real
management actions, and the real action-effect/cost tables in
\texttt{real\_ecology\_data}. All experiment rows in the analyzed result set ran
on the CPU/\texttt{numpy} backend.

The main takeaways are:
\begin{itemize}
    \item The real-data benchmark is implemented and has a complete first
    audited result set: 6480 main grid rows plus 1008 P5 raw-filter ablation
    rows, with 0 failures.
    \item The safe-mode penalty is locked at \texttt{collapse\_penalty = 5} as
    an average-case operating point.
    \item At P5, the best general learner beats the best ecological baseline on
    96/112 recoverable cells, but not on the two sink populations.
    \item Learned filtering gives no material control-return gain over raw
    observations at P5; this is a control-return result, not a state-estimation
    accuracy result.
    \item A small representative P5 trajectory demo was run after the audited
    tables: two populations, four methods, five paired episodes per method.
    It provides illustrative true reward/action-over-time figures for the
    presentation, but it is not used for the headline statistical claims.
\end{itemize}

\section{Experiment Setting}

\begin{table}[htbp]
\centering
\caption{Real-ecology experiment setting.}
\begin{tabular}{ll}
\toprule
Item & Setting \\
\midrule
Run directory & \texttt{real\_ecology\_runs/psafe\_overnight\_20260705/} \\
Data source & Frozen \texttt{code/real\_ecology\_data} table \\
Populations & 9 real populations, including 2 demographic sinks \\
Actions & 11 real management actions \\
Families & Ricker, Allee, theta, regime \\
Observation noise & $\sigma_{\mathrm{obs}}\in\{0.0,0.1,0.2,0.4\}$ \\
Methods & MOPO, RefPlan, BA-MCTS, PLUS, MOOR, Delphic, OGSRL \\
Reward modes & safe; yield only at P10 for comparison \\
Safety penalty mode & occupancy \\
Penalty grid & \texttt{collapse\_penalty} $\in\{2,5,10,20\}$ \\
Filter policy for P decision & General methods learned; PLUS/MOOR ricker \\
Backend & \texttt{numpy} for all 7488 analyzed rows \\
Evaluation budget & 5 seeds $\times$ 4 episodes $\times$ horizon 50 \\
Training/data budget & 4000 transitions, episode length 25 \\
Planner budget & horizon 5, 96 sequences, 32 particles \\
\bottomrule
\end{tabular}
\end{table}

\begin{table}[htbp]
\centering
\caption{Artifact provenance and completion.}
\begin{tabular}{ll}
\toprule
Artifact & Value \\
\midrule
\texttt{all\_metrics.csv} rows & 7488 \\
Main grid rows & 6480 \\
P5 raw-filter ablation rows & 1008 \\
Failures & 0 \\
Backend labels & all \texttt{numpy} \\
Data-table labels & single frozen \texttt{real\_ecology\_data} path \\
Tests & 37/37 real-ecology tests passing \\
Aggregate path & fixed in live package; frozen snapshot kept unchanged \\
Not run & calibration gates, PLUS-GPU, sink action diagnostic, pooled agent \\
Representative trajectory demo & 8 trace rows; 2 populations $\times$ 4 methods; 5 episodes each \\
\bottomrule
\end{tabular}
\end{table}

\section{Plot Support In The Current Artifacts}

The audited full-grid artifacts support aggregate reward plots, method
comparison plots, training-convergence overview plots, final holdout
validation-error plots, and aggregate action summaries. The full-grid
\texttt{episodes.csv} files do not contain true per-timestep trajectories:
they store per-episode aggregates and danger-zone action fractions, not a full
trajectory table with timestep, state, action, and reward.

For presentation visuals, a separate small P5 trajectory demo was run after the
audited result package. It reuses the frozen code snapshot and cached P5
datasets, captures true per-timestep state/action/reward traces for Puerto
Rican parrot and Egyptian vulture under \texttt{ricker},
$\sigma_{\mathrm{obs}}=0.2$, \texttt{filter=learned}, and methods OGSRL, MOPO,
PLUS, and MOOR. These plots are illustrative trajectory examples, not
replacement evidence for the 7488-row aggregate analysis.

\begin{table}[htbp]
\centering
\caption{Which requested plots are supported by the current run.}
\begin{tabular}{p{0.28\linewidth}p{0.2\linewidth}p{0.42\linewidth}}
\toprule
Requested plot & Supported now? & What is available \\
\midrule
Reward over time & Representative only & Full grid has aggregate rewards; the P5 demo has true per-step reward traces for 2 populations $\times$ 4 methods. \\
Actions over time & Representative only & Full grid has action entropy and danger-zone action fractions; the P5 demo has per-step action traces and action-composition plots. \\
Training convergence & Partly & Overview training-loss plots exist for penalties P2/P5/P10/P20; multi-step histories are method-dependent. \\
Validation error & Partly & Final holdout dynamics RMSE/log-MSE are available; full validation curves are not universal. \\
\bottomrule
\end{tabular}
\end{table}

For a future full-grid trajectory analysis, the main evaluator should log
per-step records directly. The present trajectory figures are produced by the
targeted trace-capture scripts in the trajectory-demo directory.

\section{P\_safe Penalty Decision}

Table~\ref{tab:psafe} reports the filter-consistent decision subset: general
learners use \texttt{filter=learned}; PLUS and MOOR use \texttt{filter=ricker};
the P5 raw rows are excluded from penalty selection.

\begin{longtable}{llrrrrrr}
\caption{P\_safe decision populations in safe mode.}\label{tab:psafe}\\
\toprule
Population & P & Return & Collapse & Unsafe & Final & Min & Cost \\
\midrule
\endfirsthead
\toprule
Population & P & Return & Collapse & Unsafe & Final & Min & Cost \\
\midrule
\endhead
Amur tiger & 2 & 2.84 & 0.095 & 0.021 & 83.7 & 56.8 & 6.842 \\
Amur tiger & 5 & \textbf{2.17} & 0.057 & 0.018 & 97.1 & 64.4 & 7.911 \\
Amur tiger & 10 & 0.85 & 0.053 & 0.019 & 97.2 & 65.6 & 8.462 \\
Amur tiger & 20 & -2.51 & 0.065 & 0.024 & 105.4 & 67.4 & 9.079 \\
Puerto Rican parrot & 2 & 2.57 & 0.176 & 0.070 & 158.5 & 65.0 & 3.443 \\
Puerto Rican parrot & 5 & \textbf{0.62} & 0.163 & 0.065 & 162.5 & 66.2 & 3.987 \\
Puerto Rican parrot & 10 & -2.29 & 0.161 & 0.063 & 181.3 & 66.7 & 4.897 \\
Puerto Rican parrot & 20 & -8.25 & 0.159 & 0.062 & 182.8 & 66.8 & 5.393 \\
Egyptian vulture & 2 & -37.63 & 0.000 & 1.000 & 2.8 & 2.1 & 3.928 \\
Egyptian vulture & 5 & -93.35 & 0.000 & 1.000 & 3.0 & 2.1 & 4.993 \\
Egyptian vulture & 10 & -186.04 & 0.000 & 1.000 & 2.8 & 2.1 & 6.133 \\
Egyptian vulture & 20 & -370.91 & 0.000 & 1.000 & 3.2 & 2.6 & 7.163 \\
\bottomrule
\end{longtable}

The decision rule is the largest tested P that keeps both recoverable
decision-population mean safe returns non-negative while reducing collapse
entries versus P2. This selects \texttt{collapse\_penalty = 5}. The choice is
average-case: individual cells remain negative, especially for the Puerto Rican
parrot.

\begin{figure}[htbp]
\centering
\includegraphics[width=0.82\linewidth]{report_figures/fig_psafe_decision_returns.png}
\caption{Safe return across the penalty grid for the three decision populations.}
\label{fig:psafe-returns}
\end{figure}

\section{General Learners Versus Ecological Baselines}

At P5 in safe mode, the comparison is made on matched
population/family/noise cells. General learners use learned filtering; PLUS and
MOOR use the ricker filter to represent the ecological baselines.

\begin{table}[htbp]
\centering
\caption{Best general learner versus best ecological baseline at P5.}
\label{tab:family}
\begin{tabular}{lrrrrr}
\toprule
Scope & Cells & Best general & Best baseline & $\Delta$ & Wins \\
\midrule
All 9 populations & 144 & -3.93 & -4.96 & +1.04 & 102/144 (70.8\%) \\
Recoverable only & 112 & 7.74 & 6.28 & +1.46 & \textbf{96/112 (85.7\%)} \\
Sinks only & 32 & -44.76 & -44.30 & -0.46 & 6/32 (18.8\%) \\
\bottomrule
\end{tabular}
\end{table}

\begin{table}[htbp]
\centering
\caption{Per-method P5 safe comparison against PLUS-ricker and MOOR-ricker.}
\label{tab:methods}
\begin{tabular}{lrrrrrr}
\toprule
Method & Mean return & $\Delta$ PLUS & $\Delta$ MOOR & Beats both & Recov. beats & Collapse \\
\midrule
OGSRL & -5.06 & +0.16 & +0.26 & 81/144 & 79/112 & 0.014 \\
BA-MCTS & -5.69 & -0.47 & -0.37 & 68/144 & 64/112 & 0.047 \\
RefPlan & -5.31 & -0.08 & +0.02 & 38/144 & 37/112 & 0.034 \\
MOPO & -5.45 & -0.22 & -0.13 & 34/144 & 34/112 & 0.036 \\
Delphic & -9.77 & -4.54 & -4.44 & 33/144 & 33/112 & 0.188 \\
\bottomrule
\end{tabular}
\end{table}

\begin{figure}[htbp]
\centering
\includegraphics[width=0.9\linewidth]{report_figures/fig_p5_return_by_method.png}
\caption{P5 safe return by method and population scope. Recoverable populations
show positive returns; sinks dominate the negative all-population averages.}
\label{fig:p5-return-method}
\end{figure}

\section{Learned-Versus-Raw Filter Ablation}

The P5 raw-filter ablation compares learned filtering to raw observations on the
same population/family/noise/method grid. The result is a clean null result for
control return.

\begin{table}[htbp]
\centering
\caption{P5 learned-vs-raw control ablation.}
\label{tab:raw}
\begin{tabular}{lrrrr}
\toprule
Method & Learned & Raw & Learned--raw & Unsafe learned/raw \\
\midrule
BA-MCTS & -5.69 & -5.37 & -0.32 & 0.126 / 0.126 \\
Delphic & -9.76 & -9.10 & -0.67 & 0.185 / 0.199 \\
MOOR & -4.50 & -4.69 & +0.18 & 0.112 / 0.115 \\
MOPO & -5.45 & -5.42 & -0.03 & 0.121 / 0.122 \\
OGSRL & -5.06 & -4.49 & -0.58 & 0.113 / 0.112 \\
PLUS & -5.23 & -5.73 & +0.50 & 0.129 / 0.136 \\
RefPlan & -5.31 & -5.29 & -0.02 & 0.121 / 0.121 \\
\bottomrule
\end{tabular}
\end{table}

\begin{figure}[htbp]
\centering
\includegraphics[width=0.72\linewidth]{report_figures/fig_learned_vs_raw_sigma.png}
\caption{Learned-minus-raw return by observation-noise level at P5. The
difference stays small even at $\sigma_{\mathrm{obs}}=0.4$.}
\label{fig:raw-sigma}
\end{figure}

\section{Sink Populations}

The two sink populations are reported separately and are not pooled with
recoverable populations.

\begin{table}[htbp]
\centering
\caption{P5 safe sink metrics under the filter-consistent policy subset.}
\label{tab:sinks}
\begin{tabular}{lrrrrp{0.36\linewidth}}
\toprule
Sink & Unsafe & Min & Final & Return & Regime \\
\midrule
Egyptian vulture & 1.000 & 2.1 & 3.0 & -93.3 & Always below the safety floor; penalty-dominated. \\
Bottlenose dolphin & 0.097 & 10.4 & 13.2 & -2.8 & Mostly above relative safety floor, below absolute MVP floor. \\
\bottomrule
\end{tabular}
\end{table}

\section{Action Summaries}

The audited full-grid run records aggregate danger-zone action fractions in
each \texttt{episodes.csv}. Figure~\ref{fig:action-profile} shows that
full-grid action summary at P5. The representative trajectory demo in the next
section adds true per-timestep action traces for two selected cells.

\begin{figure}[htbp]
\centering
\includegraphics[width=0.88\linewidth]{report_figures/fig_p5_danger_action_profile.png}
\caption{P5 aggregate danger-zone action fractions by method. This is not an
action-over-time plot; it is the action summary currently supported by the
logged artifacts.}
\label{fig:action-profile}
\end{figure}

\section{Representative Reward And Action Trajectories}

Figures~\ref{fig:parrot-trajectory}--\ref{fig:vulture-actions} come from the
small P5 trajectory demo, not the full aggregate sweep. The demo uses
Puerto Rican parrot as a recoverable near-margin population and Egyptian
vulture as a sink population. Each plotted method is refit on the same frozen
P5 cached dataset used by the audited run, then rolled for five paired
evaluation episodes. Lines show medians over those paired episodes; bands show
the interquartile range where applicable.

\begin{figure}[htbp]
\centering
\includegraphics[width=0.94\linewidth]{report_figures/puerto_rican_parrot_ricker_sigma_0p2_trajectory.png}
\caption{Representative P5 Puerto Rican parrot trajectory demo under
\texttt{ricker}, $\sigma_{\mathrm{obs}}=0.2$, and learned filtering. Panels show
true state, true immediate reward, cumulative true reward, and median action id
over time for OGSRL, MOPO, PLUS, and MOOR.}
\label{fig:parrot-trajectory}
\end{figure}

\begin{figure}[htbp]
\centering
\includegraphics[width=0.94\linewidth]{report_figures/puerto_rican_parrot_ricker_sigma_0p2_actions.png}
\caption{Representative P5 Puerto Rican parrot action composition over time.
Each panel gives the per-timestep fraction of the five paired demo episodes
selecting each real management action.}
\label{fig:parrot-actions}
\end{figure}

\begin{figure}[htbp]
\centering
\includegraphics[width=0.94\linewidth]{report_figures/egyptian_vulture_ricker_sigma_0p2_trajectory.png}
\caption{Representative P5 Egyptian vulture trajectory demo under
\texttt{ricker}, $\sigma_{\mathrm{obs}}=0.2$, and learned filtering. This is the
sink-population contrast to Figure~\ref{fig:parrot-trajectory}; it should not
be pooled with recoverable populations.}
\label{fig:vulture-trajectory}
\end{figure}

\begin{figure}[htbp]
\centering
\includegraphics[width=0.94\linewidth]{report_figures/egyptian_vulture_ricker_sigma_0p2_actions.png}
\caption{Representative P5 Egyptian vulture action composition over time.
The sink case is shown separately because its safety and reward regime differs
qualitatively from recoverable populations.}
\label{fig:vulture-actions}
\end{figure}

\section{Training Convergence And Validation Error}

The run contains overview convergence plots for each penalty and per-cell
\texttt{training\_history.png} files for learned/model cells. The convergence
overview is now limited to the two iteratively trained methods with meaningful
step-wise signals: OGSRL uses training return, where higher is better, and
Delphic uses Bellman MSE, where lower is better. Closed-form or one-shot
fitters (MOOR, MOPO, RefPlan, BA-MCTS, and PLUS) do not have a descent curve;
their fit quality is summarized by the holdout-RMSE figure.

\begin{figure}[htbp]
\centering
\includegraphics[width=0.94\linewidth]{report_figures/convergence_p2.png}
\includegraphics[width=0.94\linewidth]{report_figures/convergence_p5.png}
\caption{Training convergence for the two iteratively trained methods, shown as
median plus interquartile range over cells. OGSRL plots training return
(higher is better); Delphic plots Bellman MSE (lower is better). Numerically
unstable Delphic cells are disclosed in each panel. This figure shows P2 and
P5; Figure~\ref{fig:convergence-highp} shows P10 and P20. Closed-form and
one-shot fitters do not have descent curves; see
Figure~\ref{fig:holdout-rmse} for fit quality.}
\label{fig:convergence}
\end{figure}

\begin{figure}[htbp]
\centering
\includegraphics[width=0.94\linewidth]{report_figures/convergence_p10.png}
\includegraphics[width=0.94\linewidth]{report_figures/convergence_p20.png}
\caption{Training convergence for P10 and P20, using the same median/IQR
format as Figure~\ref{fig:convergence}. Delphic instability counts are 48/288
cells at P10 and 28/144 cells at P20; the P2/P5 counts are 23/144 and 23/288.}
\label{fig:convergence-highp}
\end{figure}

\begin{figure}[htbp]
\centering
\includegraphics[width=0.76\linewidth]{report_figures/fig_p5_holdout_rmse_by_method.png}
\caption{P5 final holdout dynamics RMSE by method. This is a final validation
error summary, not a full validation-over-time curve for every method.}
\label{fig:holdout-rmse}
\end{figure}

\section{Potential Improvements}

The current experiment is already sufficient to demonstrate the move from the
dummy setting to real ecology data, but several natural improvements remain.
They fall into two groups: adapting algorithms more directly to the new setting,
and broadening the experiment design.

\subsection{Algorithm Adaptation}

\begin{itemize}
    \item \textbf{Safety-aware learning objective.} The P5 result shows that
    safe-mode performance is sensitive to the occupancy penalty and to
    near-threshold populations such as the Puerto Rican parrot. Future learners
    could optimize a constrained or risk-sensitive objective directly, rather
    than treating the safety term as an ordinary scalar reward penalty.
    \item \textbf{Sink-population policy adaptation.} The Egyptian vulture and
    bottlenose dolphin behave differently from recoverable populations. A
    dedicated sink diagnostic, with whole-episode action logging, should test
    whether policies learn appropriate rescue/translocation behavior or simply
    become penalty-dominated.
    \item \textbf{Real-action structure.} The 11 real management actions have
    semantic groups: no action, harvest/revenue, conservation support, habitat
    or capacity improvement, and translocation. Algorithms could exploit this
    structure through action embeddings, hierarchical action selection, or
    cost-aware policy heads instead of treating action ids as unrelated labels.
    \item \textbf{Population-conditioned or pooled learning.} This run uses
    one agent per cell. A future pooled or multi-task agent could share
    statistical strength across populations while conditioning on
    population-specific quantities such as $K$, $N_0$, safety floor, and
    action-effect caps. This should be introduced after the per-cell baseline
    is stable.
    \item \textbf{Estimation-control separation.} The learned-vs-raw ablation
    found no control-return gain from learned filtering. A follow-up should
    evaluate state-estimation accuracy directly, not only control return, so
    any filter-value claim is tied to the correct metric.
\end{itemize}

\subsection{Experiment Setting}

\begin{itemize}
    \item \textbf{Calibration gates at P5.} The selected P5 penalty is an
    average-case choice. Additional calibration gates could report whether it
    is also acceptable under stricter worst-case or population-specific
    criteria.
    \item \textbf{Per-step logging as standard output.} The representative
    trajectory demo required a separate trace-capture rerun. Future full runs
    should log timestep, true state, observation, belief mean, action, immediate
    true reward, and safety indicators directly, enabling reward/action-over-time
    plots without a separate script.
    \item \textbf{More representative trajectory cells.} The current trajectory
    figures cover one recoverable near-margin population and one sink. A small
    fixed presentation panel could add Amur tiger and bottlenose dolphin to
    show a healthier recoverable case and the second sink regime.
    \item \textbf{Sensitivity analysis.} The current result varies population,
    family, observation noise, method, filter, and P value. Future runs could
    add sensitivity to dataset size, process noise, horizon, planner budget,
    safety-fraction convention, and action-cost scaling.
    \item \textbf{Backend and runtime comparison.} CPU/\texttt{numpy} was fast
    enough for the audited runs. PLUS-GPU and broader backend comparisons are
    still useful as engineering diagnostics, but they are not needed for the
    present scientific claim.
\end{itemize}

\section{Conclusion}

The current result set is sufficient to present the transition from a dummy
ecology benchmark to a real-ecology-data benchmark. It supports audited tables,
aggregate reward plots, action-summary plots, convergence overview plots, and
final validation-error plots. A small targeted P5 trace rerun additionally
supports representative true reward/action-over-time visuals for two selected
populations. These trajectory figures should be presented as illustrative
examples; the quantitative claims remain the audited aggregate tables.

\end{document}
